Für die Erstellung von Verzeichnissen hält \LaTeX\ die Standardbefehle gemäß
Tabelle~\ref{tab:final:toclists} bereit. Mehr ist für den Autor nicht zu tun.
Die Einordnung der Verzeichnisse in die Gliederung der Arbeit wird im
Abschnitt~\ref{sec:content:section} behandelt.

\begin{table}
\centering
\begin{tabular}{ll}
\verb+\tableofcentents+   & Inhaltsverzeichnis           \\
\verb+\listoffigures+     & Abbildungsverzeichnis        \\
\verb+\listoftables+      & Tabellenverzeichnis          \\
\verb+\lstlistoflistings+ & Verzeichnis aller Quellcodes \\
\end{tabular}
\caption{Verzeichnisbefehle}
\label{tab:final:toclists}
\end{table}

Für Definitionen, Sätze oder selbstdefinierte Gleitumgebungen existieren
weitere Befehle. Hierzu sei auf die Dokumentation der entsprechenden Pakete
verwiesen.\todo[inline]{Pakete heraussuchen und verlinken.}

Normalerweise erscheinen nur nummerierte Abschnitt im Inhaltsverzeichnis.
Soll ein unnummerierter Punkt dennoch ins Inhaltsverzeichnis aufgenommen
werden, kann dies über den \verb+\addcontentsline+-Befehl erfolgen, wie in
Listing~\ref{lst:final:addcontentsline} gezeigt. Der Befehl muss immer nach dem
entsprechenden Gliederungsbefehl aufgerufen werden. Sonst stimmen die
Seitenverweise nicht.

\lstinputlisting[language={[LaTeX]{TeX}},style=block,escapechar=\!,float,caption={Aufnahme zusätzlicher Verzeichniseinträge},label={lst:final:addcontentsline},morekeywords={chapter,subsection}]{code/final_addcontentsline.de.tex}
